\section*{\LARGE Running the Adventure}\stepcounter{section}\phantomsection\addcontentsline{toc}{section}{Running the Adventure}
{\entryfont The following guidance is intended to help the DM keep this one-shot focused, flexible, and suitable for a single-session format. As the adventure follows a loose structure built around two fixed narrative moments and a flexible series of encounters, individual scenes can be shortened, expanded, rearranged, or adjusted depending on the pace and size of the group.}
\subsection*{Adventure Structure}\stepcounter{subsection}\phantomsection\addcontentsline{toc}{subsection}{Adventure Structure}
{\entryfont The adventure begins in Ty Coeden, where the characters take part in a gathering concerning the future of the Caelthir and are encouraged to venture beyond the village and explore what these changing circumstances might mean for their people.

Once the characters leave Ty Coeden, the adventure becomes deliberately open. The encounters involving the Human Lumber Camp, the Hunters near Loch Leven, and the Deserted Hag's Nest do not build upon one another and may occur in any order. The DM should choose or arrange them according to the characters' route through the Woods of Lomond, allowing each encounter to arise naturally along the characters' journey.

None of these encounters is required to progress the adventure. Depending on the available session time, the DM may run all three, select only those most appropriate to the characters' journey, or omit an encounter entirely. Their purpose is to explore different aspects of life beyond Ty Coeden rather than form a continuous chain of events.

After the characters have spent sufficient time exploring the surrounding region, they return to Ty Coeden for the Festival of New Roots, which serves as the narrative conclusion of the one-shot. A short cliffhanger following the festival then reveals that an unseen threat has begun to turn its attention toward the village.}
\subsection*{Pacing and Player Agency}\stepcounter{subsection}\phantomsection\addcontentsline{toc}{subsection}{Pacing and Player Agency}
{\entryfont "Stranded in Fife" is intended to give the players considerable freedom in how they approach both the changing circumstances of Ty Coeden and the people they encounter beyond it. Discussions concerning the future of the Caelthir should not have a single correct answer, and the characters' opinions may differ considerably from one another. Allow these disagreements to develop naturally without forcing the group toward a predetermined conclusion.

Social encounters should likewise be resolved primarily through the characters' actions and arguments rather than a single ability check. Checks can support an approach, reveal additional information, or influence uncertain NPCs, but should rarely replace roleplaying entirely.

As the adventure is intended for a single session, the DM should remain mindful of pacing. Scenes that generate strong player engagement may be allowed additional time, while travel, investigation, or optional encounters can be shortened or omitted as necessary. Enough time should always remain for the characters' return to Ty Coeden and the Festival of New Roots, as these provide the narrative conclusion of the one-shot.}
\subsection*{Characters and Encounters}\stepcounter{subsection}\phantomsection\addcontentsline{toc}{subsection}{Characters and Encounters}
{\entryfont "Stranded in Fife" is designed to be played using the pregenerated inhabitants of Ty Coeden provided with this adventure. Each player chooses one of these characters for the duration of the one-shot instead of using their regular campaign character. As the available characters represent different backgrounds, personalities, and relationships with Ty Coeden, their individual perspectives can strongly influence how they approach the questions raised throughout the adventure.

No particular pregenerated character is required for any scene to function. The adventure deliberately avoids relying on a specific class feature, spell, proficiency, or character being present, allowing the players to freely choose whichever characters appeal to them. The DM should likewise avoid assuming that any particular character accompanies the group when presenting information or resolving encounters.

Encounters throughout this adventure place greater emphasis on interaction, investigation, and creative problem-solving than on combat. Many situations have no predetermined solution and can be approached through conversation, skill checks, magic, or unconventional ideas. Combat is used sparingly and is intended primarily to provide variety rather than form the central challenge of the adventure.}
\subsection*{Foreshadowing the Zircon Foxes}\stepcounter{subsection}\phantomsection\addcontentsline{toc}{subsection}{Foreshadowing the Zircon Foxes}
{\entryfont The surviving Cobalt Fox pups provide a subtle background thread throughout Stranded in Fife, but they are not intended to become an active threat during the adventure itself. Their presence should be hinted at only indirectly through unusual tracks, disturbed wildlife, or other signs encountered within the Woods of Lomond and around Loch Leven. These clues should create unease without allowing the characters to determine exactly what is responsible.

Even particularly successful attempts to investigate or follow these signs should reveal additional details rather than lead directly to the creatures themselves. The Zircon Foxes remain beyond the characters' reach during the one-shot, and the DM should avoid providing a clear description or identification of them. Likewise, not every unusual occurrence within the woods should be connected to the foxes; their presence works best when it remains only one subtle element among the wider events of the adventure.

The first unmistakable indication of what has emerged from the surviving pups is reserved for the closing cliffhanger. Until that moment, the characters may suspect that something unusual is moving through the forest, but neither they nor the inhabitants of Ty Coeden should understand the danger that is beginning to develop.}